\section{CD Repo Và GitOps Desired State}

CD repo \texttt{emanhthangngot/yas-cd} là source of truth cho cluster. Jenkins không
deploy trực tiếp vào namespace \texttt{dev}, \texttt{staging}, \texttt{developer};
Jenkins chỉ commit thay đổi vào CD repo, sau đó ArgoCD reconcile.

\subsection{Cấu Trúc Desired State}

\begin{table}[H]
\centering
\caption{Vai trò các thư mục GitOps chính}
\begin{tabular}{p{4cm} p{10cm}}
\hline
\textbf{Đường dẫn} & \textbf{Vai trò} \\
\hline
\path{base/} & Config dùng chung cho application workloads: route config, compatibility service, shared config. \\
\hline
\path{platform/base/} & PostgreSQL, Redis, Kafka, Elasticsearch, Keycloak, ingress và namespace platform. \\
\hline
\path{overlays/dev} & Môi trường dev, auto-sync, nhận image tag từ main build. \\
\hline
\path{overlays/staging} & Môi trường staging, manual sync, chỉ dùng immutable release tag. \\
\hline
\path{overlays/developer} & Developer preview, dormant mặc định, active khi chạy \texttt{developer\_build}. \\
\hline
\path{argocd/apps} & ArgoCD Application manifests trỏ về CD repo \texttt{main}. \\
\hline
\path{operations/sampledata-seed} & Job seed dữ liệu một lần cho dev/staging. \\
\hline
\end{tabular}
\end{table}

\subsection{CQ Service Scope}

Theo file PDF phân loại service CQ, các service cốt lõi gồm product, cart, order,
customer, inventory, tax, media, search, storefront/backoffice BFF và UI, swagger UI,
sampledata. CD repo giữ thêm \texttt{location}, \texttt{payment},
\texttt{payment-paypal} vì chúng là runtime dependency của tax/order flow.

Các service optional như \texttt{promotion}, \texttt{rating}, \texttt{recommendation}
và \texttt{webhook} không chạy trong baseline để giảm tải single-node, nhưng vẫn còn
trong source và chart để mở rộng về sau.

\subsection{Validation Gates Trong CD Repo}

Trước khi commit GitOps, nhóm dùng các script:

\begin{itemize}\sloppy
    \item \texttt{scripts/}\allowbreak\texttt{validate-gitops.sh}: render overlays, kiểm tra service catalog, secret-like patterns, sampledata operation overlays và active environment policy.
    \item \texttt{scripts/}\allowbreak\texttt{validate-}\allowbreak\texttt{staging-}\allowbreak\texttt{immutable.sh}: đảm bảo staging chỉ dùng release tag immutable, không dùng \texttt{latest}.
    \item \texttt{scripts/}\allowbreak\texttt{update-image-tag.sh}: contract cho Jenkins main/dev update.
    \item \texttt{scripts/}\allowbreak\texttt{promote-staging-release.sh}: đổi toàn bộ staging image tag sang \texttt{vX.Y.Z}.
\end{itemize}

\lstinputlisting[
  caption={GitOps validation passed trước khi commit desired state},
  basicstyle=\ttfamily\scriptsize
]{../evidence/07_validate_gitops_passed.txt}

\subsection{ArgoCD Sync Policy}

\begin{itemize}
    \item \texttt{yas-dev}: automated sync/self-heal cho main branch changes.
    \item \texttt{yas-staging}: manual sync để operator approve release.
    \item \texttt{yas-developer}: automated sync nhưng dormant mặc định.
    \item \texttt{yas-platform}: quản lý shared infra.
\end{itemize}

\begin{figure}[H]
\centering
\reportimg{img/argocd/24_argocd_outofsync_diff.png}
\caption{Staging release tạo diff trong ArgoCD trước khi operator sync thủ công}
\end{figure}
